\exheading{17}{The Choice of the Damping Factor $m$}{Ömer Özkan}

\exblock{Problem Statement}

We are asked how the value of $m$ should be chosen, and how this choice affects the rankings and the computation time.
We answer with an experiment (script \texttt{ex17.py}): we rank two webs for
$m \in \{0.05,\ 0.15,\ 0.25,\ 0.50,\ 0.85,\ 0.99\}$,
\begin{enumerate}
  \item the 4-page web of Figure~2.1,
  \item the Hollins dataset (6012 pages),
\end{enumerate}
and we record the number of iterations, the time and the ranking.
The power method starts from the uniform vector and stops when $\|\vx_{k+1}-\vx_k\|_1 < \texttt{tol} = 10^{-10}$.
We state the tolerance because the number of iterations has no meaning without it.

\exblock{Results: 4-Page Web (Figure 2.1)}

\subsubsection*{Convergence Speed}

Even on this tiny web the trend is clear: a larger $m$ needs fewer iterations (Table~\ref{tab:ex17-fig21}).

\begin{table}[H]\centering
\begin{tabular}{ccccccc}
\toprule
$\mathbf{m}$ & \textbf{Iterations} & \textbf{Time (s)} & $\mathbf{x_1}$ & $\mathbf{x_2}$ & $\mathbf{x_3}$ & $\mathbf{x_4}$ \\
\midrule
0.05 & 35 & 0.0006 & 0.380907 & 0.133120 & 0.289620 & 0.196353 \\
0.15 & 31 & 0.0004 & 0.368151 & 0.141809 & 0.287962 & 0.202078 \\
0.25 & 26 & 0.0003 & 0.354915 & 0.151229 & 0.285917 & 0.207940 \\
0.50 & 18 & 0.0002 & 0.320064 & 0.178344 & 0.278662 & 0.222930 \\
0.85 &  9 & 0.0002 & 0.269896 & 0.225995 & 0.261165 & 0.242944 \\
0.99 &  5 & 0.0001 & 0.251256 & 0.248338 & 0.250827 & 0.249579 \\
\bottomrule
\end{tabular}
\caption{Figure~2.1 for several values of $m$ ($\texttt{tol} = 10^{-10}$); the order is always $1>3>4>2$}
\label{tab:ex17-fig21}
\end{table}


\subsubsection*{Rankings}

The order $1 > 3 > 4 > 2$ never changes.
What changes is the distance between the scores: for $m = 0.05$ page~1 has $0.381$ and page~2 has $0.133$, for $m = 0.99$ all four scores are within $0.003$ of the uniform value $0.25$.
This is what formula (3.2) says: a part $m$ of every score is the same for all pages, so when $m$ grows the scores move towards $1/n$.

\exblock{Results: Hollins Dataset (6012 Pages)}

\subsubsection*{Convergence Speed}

On the large web the saving is important: going from $m = 0.15$ to $m = 0.85$ divides the number of iterations by almost forty, and the time goes from $2.9$~s to $0.15$~s (Table~\ref{tab:ex17-hollins}).

\begin{table}[h]\centering
\begin{tabular}{cccl}
\toprule
\textbf{m} & \textbf{Iterations} & \textbf{Time (s)} & \textbf{Top 5 pages} \\
\midrule
0.05 & 2850 & 36.33 & 4023, 3227, 4075, 5254, 3834 \\
0.15 &  407 &  2.91 & 2, 37, 38, 61, 52 \\
0.25 &  152 &  1.56 & 2, 37, 38, 61, 52 \\
0.50 &   39 &  0.45 & 2, 425, 37, 38, 52 \\
0.85 &   11 &  0.15 & 2, 425, 37, 38, 52 \\
0.99 &    5 &  0.07 & 2, 425, 37, 38, 52 \\
\bottomrule
\end{tabular}
\caption{Hollins for several values of $m$ (\texttt{tol} $= 10^{-10}$, uniform start)}
\label{tab:ex17-hollins}
\end{table}




\subsubsection*{Rankings}

The home page (page~2) is first for every $m \ge 0.15$, but the rest of the top five moves at $m = 0.50$.
Page~425 (a library page with a list of web links) is not in the top ten for $m = 0.15$ and 2nd for $m \ge 0.50$.
Looking at the data, page~425 has only 87 backlinks, but 57 of them come from pages with one or two links, so it receives their vote almost undivided.
For small $m$ these small pages have a low score and their votes are worth little; for large $m$ all scores are close to $1/n$, every voter is worth the same, and the undivided votes of page~425 push it up.
So on a real web a large $m$ does not only flatten the scores, it also changes the order.

For $m = 0.05$ the top five are pages of one course site (\texttt{faculty/saloweyca/clas 395}), which link only to each other: with dangling pages $\M$ loses vote at every step, and for very small $m$ the pages that keep their vote come first.

\exblock{Conclusion}

A larger $m$ makes the power method faster but moves every score towards $1/n$, so the ranking loses the information of the links.
On a small web the order does not change; on Hollins the top pages change from $m = 0.50$, and for a very small $m$ the ranking is decided by the dangling pages.
The value $m = 0.15$ used by Google is a reasonable compromise: the ranking follows the link structure, and the method converges in about four hundred iterations on a web of six thousand pages.